% !TeX root = ../main.tex

\documentclass[../main.tex]{subfiles}
\begin{document}

\begin{center}
    \Large{\textbf{LỜI NÓI ĐẦU}}\\
\end{center}
\vspace{1cm}

Trong nhà máy nhiệt điện, hệ thống cung cấp dầu nhiên liệu khởi động cho lò than phun giữ vai trò đặc biệt quan trọng ở giai đoạn khởi động, ổn định cháy và chuyển đổi chế độ vận hành. Đối tượng này vừa liên quan trực tiếp đến chất lượng quá trình cháy, vừa ảnh hưởng đến mức tiêu hao nhiên liệu, độ tin cậy của thiết bị và an toàn vận hành toàn tổ máy. Do dầu nhiên liệu nặng có độ nhớt lớn, hệ thống cấp dầu phải làm việc trong điều kiện có gia nhiệt, tuần hoàn hồi lưu và kiểm soát áp suất chặt chẽ; vì vậy đây là một bài toán điều khiển quá trình điển hình trong công nghiệp nhiệt điện.

Xuất phát từ yêu cầu thực tế đó, bài tập lớn này tập trung nghiên cứu, mô hình hóa và thiết kế hệ thống điều khiển cho quá trình cấp dầu nhiên liệu khởi động của lò than phun. Trên cơ sở khảo sát công nghệ, các thông số vận hành và đặc tính động học của đối tượng, đề tài tiến hành nhận dạng mô hình bằng dữ liệu thực nghiệm xử lý trên môi trường Python, sau đó tổng hợp bộ điều khiển theo hai phương pháp là IMC và điều khiển bền vững. Việc lựa chọn hai phương pháp này nhằm bảo đảm hệ thống có khả năng đáp ứng nhanh, bền vững trước nhiễu và phù hợp với đặc tính quán tính lớn, có trễ của đối tượng nhiệt.

Trong quá trình thực hiện, nội dung bài tập lớn được triển khai theo hướng gắn kết chặt chẽ giữa công nghệ, mô hình toán và giải pháp điều khiển. Các thông số đo lường chủ yếu như lưu lượng, áp suất, nhiệt độ và mức nhiên liệu được phân tích để xác định biến điều khiển, biến nhiễu và yêu cầu bảo vệ. Từ đó, hệ thống điều khiển được đề xuất theo cấu trúc phù hợp với thực tế vận hành của nhà máy nhiệt điện, đồng thời bảo đảm tính khả thi khi triển khai trên nền tảng điều khiển công nghiệp.

Bên cạnh giá trị chuyên môn, đề tài còn giúp củng cố phương pháp tư duy kỹ thuật trong phân tích quá trình công nghệ, xây dựng mô hình và đánh giá chất lượng điều khiển. Đây là những năng lực cốt lõi đối với sinh viên kỹ thuật khi tiếp cận các hệ thống tự động hóa công nghiệp quy mô lớn. Kết quả của bài tập lớn vì vậy không chỉ có ý nghĩa học thuật mà còn có thể sử dụng trực tiếp làm tài liệu tham khảo cho các bài toán điều khiển quá trình trong thực tế.

Nội dung bài tập lớn được bố cục thành bốn chương chính:
\begin{itemize}[noitemsep,topsep=5pt]
    \item \textbf{Chương 1 --- Tổng quan về đối tượng công nghệ:} trình bày tổng quan công nghệ cấp dầu nhiên liệu khởi động cho lò than phun, các thông số cần giám sát, vai trò của hệ thống và các yếu tố ảnh hưởng đến quá trình.
    \item \textbf{Chương 2 --- Thiết kế vòng điều khiển đối tượng:} phân tích triết lý thiết kế, lựa chọn thiết bị đo và cơ cấu chấp hành, xác định dải làm việc, cảnh báo và bảo vệ, đồng thời trình bày phương pháp xử lý số liệu thực nghiệm.
    \item \textbf{Chương 3 --- Phân tích và tổng hợp hệ thống:} nhận dạng đối tượng, xây dựng mô hình toán học, phân tích cấu trúc điều khiển, tổng hợp bộ điều khiển theo IMC và điều khiển bền vững, đánh giá chất lượng hệ thống bằng mô phỏng Python.
    \item \textbf{Chương 4 --- Kết luận và kiến nghị:} tổng hợp kết quả đạt được, đánh giá ý nghĩa kỹ thuật và đề xuất hướng phát triển tiếp theo.
\end{itemize}

\vspace{1cm}

\end{document}